% !TeX program = xelatex
% !TeX encoding = UTF-8
% !TeX spellcheck = en_US
% !BIB program = biber
%% 
%% The above lines help editors like TeXstudio to automatically choose the right tools
%% to compile your LaTeX source file. If your tool does not support these magic comments,
%% you will need to make appropriate manual choices.
%% 
%% You can safely use "pdflatex" instead of "xelatex" if you prefer the pdfLaTeX toolchain.
%% However, pdfLaTeX will not be able to deliver the professional font experience that you
%% will get with XeLaTeX. You can also safely use "lualatex" instead of "xelatex" while
%% preserving the professional font experience if you prefer the LuaLaTeX toolchain.
%% 
%% _Important_: These magic comments should be on the first lines of your source file.
%% 
%%%%%%%%%%%%%%%%%%%%%%%%%%%%%%%%%%%%%%%%%%%%%%%%%%%%%%%%%%%%%%%%%%%%%%%%%%%%%%%%

%%%%%%%%%%%%%%%%%%%%%%%%%%%%%%%%%%%%%%%%%%%%%%%%%%%%%%%%%%%%%%%%%%%%%%%%%%%%%%%%
%% 
%%            JJJJ   K                         K   UUUU         UUUU  
%%            JJJJ   KKKK                   KKKK   UUUU         UUUU  
%%            JJJJ   KKKKKK               KKKKKK   UUUU         UUUU  
%%            JJJJ      KKKKKK         KKKKKK      UUUU         UUUU  
%%            JJJJ         KKKKKK   KKKKKK         UUUU         UUUU  
%%            JJJJ            KKKKKKKKK            UUUU         UUUU  
%%    JJ     JJJJJ               KKK               UUUUU       UUUUU  
%%  JJJJJJJJJJJJJ    KKKKKKKKKKKKKKKKKKKKKKKKKKK    UUUUUUUUUUUUUUU   
%%    JJJJJJJJJ      KKKKKKKKKKKKKKKKKKKKKKKKKKK      UUUUUUUUUUU     
%% 
%% This is an example file for using the JKU LaTeX technical report template
%% for your seminar report.
%% 
%% Template created by Michael Roland (2023)
%% 
%%%%%%%%%%%%%%%%%%%%%%%%%%%%%%%%%%%%%%%%%%%%%%%%%%%%%%%%%%%%%%%%%%%%%%%%%%%%%%%%

%%%%%%%%%%%%%%%%%%%%%%%%%%%%%%%%%%%%%%%%%%%%%%%%%%%%%%%%%%%%%%%%%%%%%%%%%%%%%%%%
%% 
%% Document class: This is a koma-script article.
%% 
\documentclass[a4paper,oneside,10pt,ngerman,english]{scrartcl}
%% 
%% The comma-separated list in square brackets are class options.
%% Useful options that you might want to use:
%% 
%% Paper size:
%%  * a4paper ... A4 paper size
%% 
%% Optimize for single-sided or double-sided printing:
%%  * oneside ... single-sided
%%  * twoside ... double-sided
%% 
%% Base font size:
%%  * 10pt ... 10-pt font is used for normal text
%%  * 11pt ... 11-pt font is used for normal text
%% 
%% Define document languages (the last specified language becomes the document default
%% language):
%%  * ngerman ... German
%%  * english ... English
%%  * ...
%% 
%% Alternate document classes: The JKU report template supports the koma-script classes
%% `scrartcl', `scrreprt' and `scrbook'. The article class `scrartcl' is well-suited
%% for a typical seminar report. However, `scrbook' or `scrreprt' may be better
%% suited for longer reports since they permit structuring your work in chapters.
%%  
%% _Important_: The document class should be the first line of LaTeX code in your main
%% source file. Do not place anything but comments / magic comments above that line (unless
%% you really know what you are doing).
%% 
%%%%%%%%%%%%%%%%%%%%%%%%%%%%%%%%%%%%%%%%%%%%%%%%%%%%%%%%%%%%%%%%%%%%%%%%%%%%%%%%

%%%%%%%%%%%%%%%%%%%%%%%%%%%%%%%%%%%%%%%%%%%%%%%%%%%%%%%%%%%%%%%%%%%%%%%%%%%%%%%%
%% 
%% Treat input files as UTF-8 encoded. Make sure to always load this when you use pdfLaTeX
%% so that pdfLaTeX knows how to read and interpret characters in this source file.
%% 
\usepackage[utf8]{inputenc}
%% 
%%%%%%%%%%%%%%%%%%%%%%%%%%%%%%%%%%%%%%%%%%%%%%%%%%%%%%%%%%%%%%%%%%%%%%%%%%%%%%%%

%%%%%%%%%%%%%%%%%%%%%%%%%%%%%%%%%%%%%%%%%%%%%%%%%%%%%%%%%%%%%%%%%%%%%%%%%%%%%%%%
%% 
%% Use the JKU LaTeX technical report template for this document.
%% 
\usepackage[seminarreport,fancyfonts,TNF]{jkureport}
%% 
%% The comma-separated list in square brackets are theme options. Useful options that you
%% might want to use:
%% 
%% Document type:
%%  * phdthesis     ... PhD thesis.
%%  * mathesis      ... Master's thesis.
%%  * diplomathesis ... Diploma thesis.
%%  * bathesis      ... Bachelor's thesis.
%%  * seminarreport ... Seminar report.
%%  * techreport    ... Technical report.
%% 
%% Color scheme selection options:
%%  * JKU  ... Use JKU (gray) color scheme (this is the default if no scheme is selected).
%%  * BUS  ... Use Business School color scheme.
%%  * LIT  ... Use Linz Institute of Technology color scheme.
%%  * MED  ... Use MED faculty color scheme.
%%  * RE   ... Use RE faculty color scheme.
%%  * SOE  ... Use School of Education color scheme.
%%  * SOWI ... Use SOWI faculty color scheme.
%%  * TNF  ... Use TNF faculty color scheme.
%% 
%% Space-efficient monospace font options (requires XeTeX/LuaTeX):
%%  * compactmono   ... Use condensed fixed-width font everywhere.
%%  * nocompactverb ... Do not use condensed fixed-width font for verbatim and listings.
%% 
%% Style-breaking options:
%%  * noimprint      ... Do not insert imprint on title pages.
%%  * nojkulogo      ... Do not insert JKU & K logos on title pages.
%%  * capstitle      ... Set document title in capital letters.
%%  * nofancyfonts   ... Do not use custom TTF fonts with XeTeX/LuaTeX / supress pdfLaTeX warning.
%%  * equalmargins   ... Decrease the outer page margin to have both page margins of equal size
%%                       (the additional outer margin is intentional and to be used for
%%                       anotations; equalmargins also causes the text width to be
%%                       significantly larger than optimal for reading).
%% 
%% Experimental options:
%%  * mathastext ... Use standard document fonts (enhanced with symbols from Fira Math font
%%                   when using XeTeX/LuaTeX) in math mode.
%% 
%% Advanced options:
%%  * noautopdfinfo     ... Do not automatically try to add pdfinfo with hyperref from document
%%                          metadata fields.
%%  * logopath={<path>} ... Set the path where the theme can find its own logo resources. This
%%                          should typically be a relative path and the default is `./logos'.
%%  * fontpath={<path>} ... Set the path where the theme can find its own font resources. This
%%                          should typically be a relative path and the default is `./fonts'.
%% 
%% Hint: Boolean options can be used in the forms `option' or `option=true' the enable the
%% option and `nooption' or `option=false' to disable the option.
%% 
%%%%%%%%%%%%%%%%%%%%%%%%%%%%%%%%%%%%%%%%%%%%%%%%%%%%%%%%%%%%%%%%%%%%%%%%%%%%%%%%

%%%%%%%%%%%%%%%%%%%%%%%%%%%%%%%%%%%%%%%%%%%%%%%%%%%%%%%%%%%%%%%%%%%%%%%%%%%%%%%%
%% 
%% This is the place where you can load additional packages. If you want to load
%% a package `biblatex', you would use the command `\usepackage{biblatex}'.
%% 

\usepackage{csquotes}
\usepackage[backend=biber,citestyle=numeric,sortcites=true,maxcitenames=2,style=ACM-Reference-Format]{biblatex}
\setcounter{biburlnumpenalty}{100} %% reducing biburl* penalties typically improves URL placement in bibliography
\setcounter{biburllcpenalty}{100}
\setcounter{biburlucpenalty}{100}
\usepackage{todonotes}
\usepackage{import}
\usepackage{amsfonts}
\usepackage{subfigure}
%\usepackage{acronym}

%% 
%%%%%%%%%%%%%%%%%%%%%%%%%%%%%%%%%%%%%%%%%%%%%%%%%%%%%%%%%%%%%%%%%%%%%%%%%%%%%%%%

%%%%%%%%%%%%%%%%%%%%%%%%%%%%%%%%%%%%%%%%%%%%%%%%%%%%%%%%%%%%%%%%%%%%%%%%%%%%%%%%
%% 
%% Bibliography data files.
%% 

\addbibresource{references.bib}

%% 
%%%%%%%%%%%%%%%%%%%%%%%%%%%%%%%%%%%%%%%%%%%%%%%%%%%%%%%%%%%%%%%%%%%%%%%%%%%%%%%%

\begin{document}
%%%%%%%%%%%%%%%%%%%%%%%%%%%%%%%%%%%%%%%%%%%%%%%%%%%%%%%%%%%%%%%%%%%%%%%%%%%%%%%%
%% 
%% Report information and title page
%% 

%% Command \title{title}: sets the title of your report
\title{Watermarks for Language Models}

%% Command \titleshort{short title}: sets an abbreviated version of the report title for page heads
%\titleshort{Optional space for your abbreviated title}

%% Command \subtitle{subtitle}: sets the subtitle for seminar/technical reports (not used for theses)
\subtitle{Practical work in AI --- Documentation}

%% Command \author{name}: sets the author's name; use \prefix{} and \suffix{} to add academic titles
%%   and suffixes, use \matno{} to add the immatriculation number
\author{Tim~Michael~Toller \matno{k11842616}}

%% Command \supervisor[number,gender]{name}: sets the name of the supervisor (where number optionally
%%   defines the rank of the supervisor (1-3) and gender specifies if the supervisor is male or female
%%   to adapt gener-specific terms in German)
\supervisor[1,male]{\prefix{A.Univ.Prof., DI Dr.} Josef~Küng}

%% Command \submissiondepartment{institute or department}: set the course or department that the thesis
%%   is submitted at
\submissiondepartment{Institute for Application-oriented Knowledge Processing (FAW)}

%% Command \date{YYYY-MM-DD}: set the day of submission (defaults to today)
%\date{2020-04-09}

%% Command \abstract{text}: set the document abstract on the title page
% \abstract{Space for your (short) abstract.}

%% Command \keywords{text}: set the document keywords
%\keywords{Space for your comma-separated keywords}


%% Finally, print the title page using the above information:
\maketitle
%% 
%%%%%%%%%%%%%%%%%%%%%%%%%%%%%%%%%%%%%%%%%%%%%%%%%%%%%%%%%%%%%%%%%%%%%%%%%%%%%%%%

\cleardoubleoddpage
%% Print the table of contents
\tableofcontents

\section{Introduction}
\subsection{Motivation and Problem Statement}
\subsection{Project Goals and Scope}

\section{Background: LLM Watermarking and Exp-Min}
\subsection{What Text Watermarking Is (Generation + Detection)}
\subsection{Exp-Min Watermarking: Idea and Properties}
\subsection{Detection Principle: Test Statistic and p-value Interpretation}

\section{System Overview}
\subsection{Architecture and Components (Backend, Core, Frontend)}
\subsection{End-to-End Workflows (Generate / Detect, Streaming vs. Non-Streaming)}

\section{Backend Implementation (FastAPI)}
\subsection{Project Structure and Entry Points}
\subsection{API Design}
\subsubsection{Health and Model Info}
\subsubsection{Generation Endpoints (\texttt{/generate}, \texttt{/generate/stream})}
\subsubsection{Detection Endpoints (\texttt{/detect}, \texttt{/detect/stream})}
\subsubsection{Corruption Endpoints (\texttt{/corrupt})}
\subsection{Configuration and Runtime Settings}
\subsection{Model Backend Integration (llama.cpp / \texttt{llama-cpp-python})}
\subsection{Watermark Core Integration (Generation Path)}
\subsection{Detection Core Integration (Detection Path)}

\section{Watermarking and Detection Core}
\subsection{Watermark Parameters and their roles}
\subsection{Generation Algorithm (High-Level Pseudocode)}
\subsection{Detection Algorithm (High-Level Pseudocode)}

\section{Frontend Implementation (Vite + React)}
\subsection{UI Structure and Main Views}
\subsection{API Client and Data Flow}
\subsection{Streaming Handling and Result Visualization}
\subsection{User Inputs and Validation (Practical Defaults)}

\section{Setup and Usage}
\subsection{Requirements and Installation (Python/Node)}
\subsection{Running Backend and Frontend Locally}
\subsection{Example Requests and Expected Outputs}

\section{Limitations and Future Work}
\subsection{When Detection Fails (Short Texts, Edits, Paraphrasing)}
\subsection{Model/Backend Constraints (Need Logits, Local Performance)}
\subsection{Possible Extensions (Evaluation Suite, Robustness Tests, Multi-Model Support)}

\section{Conclusion}
\subsection{Summary of Results}
\subsection{Outlook}

%% 
%%%%%%%%%%%%%%%%%%%%%%%%%%%%%%%%%%%%%%%%%%%%%%%%%%%%%%%%%%%%%%%%%%%%%%%%%%%%%%%%

%%%%%%%%%%%%%%%%%%%%%%%%%%%%%%%%%%%%%%%%%%%%%%%%%%%%%%%%%%%%%%%%%%%%%%%%%%%%%%%%
%% 
%% Print the bibliography
%% 
%% (Optionally) let the bibliography start on a new odd page (odd is only relevant in twoside layout)
%\cleardoubleoddpage
\printbibliography
%% 
%%%%%%%%%%%%%%%%%%%%%%%%%%%%%%%%%%%%%%%%%%%%%%%%%%%%%%%%%%%%%%%%%%%%%%%%%%%%%%%%

%% Begin with the appendix part (all further sections will be appendices)
\appendix

%%%%%%%%%%%%%%%%%%%%%%%%%%%%%%%%%%%%%%%%%%%%%%%%%%%%%%%%%%%%%%%%%%%%%%%%%%%%%%%%
%% 
%% Add your appendix sections ...
%% 

%% Make sure to start the appendix on a new odd page (odd is only relevant in twoside layout)
%\cleardoubleoddpage
% \section{An Appendix}
% \label{app:an-appendix}

% Space for an appendix.
% You can have more than one appendix section.
% Appendices are, of course, optional.


% %% 
% %%%%%%%%%%%%%%%%%%%%%%%%%%%%%%%%%%%%%%%%%%%%%%%%%%%%%%%%%%%%%%%%%%%%%%%%%%%%%%%%

\cleardoubleoddpage

\end{document}
\endinput
